% Generated draft: requires mathematical and source review.
\documentclass[11pt]{article}
\usepackage[margin=25mm]{geometry}
\usepackage{fontspec}
\setmainfont{Latin Modern Roman}
\usepackage{amsmath,amssymb,mathtools,unicode-math}
\setmathfont{Latin Modern Math}
\usepackage{enumitem,needspace,xurl,hyperref}
\setlength{\parindent}{0pt}
\setlength{\parskip}{5pt}
\setlength{\emergencystretch}{3em}
\newcommand{\N}{\mathbb{N}}
\newcommand{\Z}{\mathbb{Z}}
\newcommand{\Q}{\mathbb{Q}}
\newcommand{\R}{\mathbb{R}}
\newcommand{\C}{\mathbb{C}}
\newcommand{\F}{\mathbb{F}}
\newcommand{\A}{\mathbb{A}}
\newcommand{\PP}{\mathbb{P}}
\newcommand{\E}{\mathbb{E}}
\newcommand{\eps}{\varepsilon}
\newcommand{\dd}{\,\mathrm{d}}
\begin{document}
\section*{UnsolvedMath ID: 800252}
\textbf{Source problem code:} catalogue entry
\subsection{Definitions and mathematical statement}
Let \(n\geq 1\). A function \(v:\C^n\to[-\infty,\infty)\) is plurisubharmonic if it is upper semicontinuous and its restriction to each complex affine line is subharmonic or identically \(-\infty\). Write \(\operatorname{PSH}(\C^n)\) for this class. For notation, adopt the convention \(d^c=\tfrac{i}{2}(\bar\partial-\partial)\), so that \(dd^c=i\partial\bar\partial\); the source does not specify a normalization. For smooth plurisubharmonic \(v\), the expression \((dd^c v)^n\) denotes the \(n\)-fold exterior product. For general plurisubharmonic functions, this expression requires additional hypotheses and a specified definition; it is not automatically defined.

The supplied problem asks for conditions ensuring continuity of the complex Monge--Amp\`ere operator on \(\C^n\). More explicitly, specify a class \(\mathcal A\subseteq\operatorname{PSH}(\C^n)\), a notion of convergence \(\tau\), and any additional hypotheses \(H\) such that, for every \(u\in\mathcal A\) and every sequence \((u_m)_{m\geq1}\subseteq\mathcal A\),
\[
 u_m\xrightarrow{\tau}u\quad\text{and}\quad H((u_m)_{m\geq1},u)
 \quad\Longrightarrow\quad
 (dd^c u_m)^n\longrightarrow(dd^c u)^n,
\]
with all the displayed Monge--Amp\`ere expressions well defined.

The source leaves both convergence arrows unspecified. If the intended conclusion is weak convergence of locally finite Monge--Amp\`ere measures, it can be expressed as
\[
 \lim_{m\to\infty}\int_{\C^n}\varphi\,(dd^c u_m)^n
 =\int_{\C^n}\varphi\,(dd^c u)^n
 \qquad\text{for every }\varphi\in C_c^\infty(\C^n).
\]
This is an explicit possible interpretation of the conclusion, not a convergence convention established by the supplied record. The requested hypotheses, the convergence of the functions, and the domain and interpretation of the operator remain to be specified. Thus the record determines a conditions-finding problem rather than a single fully specified implication.
\subsection{Short English statement}
Find conditions under which convergence of a sequence of plurisubharmonic functions on \(\C^n\) to a plurisubharmonic limit implies convergence of their complex Monge--Amp\`ere expressions. The source does not specify either mode of convergence or hypotheses making these expressions well defined, so these choices remain part of the formulation.
\subsection{Sources}
\begin{itemize}
\item Local UnsolvedMath registry, permanent ID \(800252\), source code \texttt{catalogue entry}, title: \emph{Several Complex Variables Problem List --- Problem 51}. The supplied record attributes the problem to Buzzard, Fornaess, and Krantz, with proposed year 1995.
\item Buzzard, Fornaess, Krantz, \emph{The Several Complex Variables Problem List} (1995), Problem 51. Supplied source URL: \url{https://arxiv.org/abs/math/9509202}. The supplied provenance records access on 2026-10-16 and extraction by \texttt{arxiv\_latex}; these are supplied provenance statements, not a claim of independent source access.
\item The local record labels the problem \texttt{open}. This is supplied registry metadata, not an independently verified assessment of the mathematical literature.
\end{itemize}
\end{document}
